\centerline{\bf \Huge Алгебраический разнобой}
\centerline{\bf \Huge }

\begin{flushright}
    \scriptsize{Обоюдна}
\end{flushright}

\problem{1!}
Петя написал на доске натуральное число А. Если его умножить на 8, то получится квадрат натурального числа. А сколько существует таких трёхзначных чисел $B$, для которых $A \cdot B$ тоже является квадратом натурального числа?

\problem{2!}
Докажите, что произведение любых четырёх последовательных целых чисел при увеличении на 1 даёт точный квадрат.

\problem{3}
Вычислите значение выражения:

(a) $(2+3)\left(2^2+3^2\right)\left(2^4+3^4\right) \ldots\left(2^{512}+3^{512}\right)$;

(б) $\dfrac{2^3-1}{2^3+1} \cdot \dfrac{3^3-1}{3^3+1} \cdot \ldots \cdot \dfrac{2023^3-1}{2023^3+1}$;

\problem{4}
Сумма трёх различных натуральных делителей нечётного натурального числа $n$ равна 10327. Какое наименьшее значение может принимать $n$ ?

\problem{5}
Верно ли, что любое натуральное число можно представить в виде $\dfrac{m n+1}{m+n}$ для некоторых натуральных чисел $m$ и $n$?
